\documentclass[11pt]{article}
\usepackage[margin=1in]{geometry}
\usepackage{amsmath,amssymb,amsthm}
\usepackage{booktabs}
\usepackage{enumitem}
\usepackage{hyperref}
\usepackage{graphicx}
\usepackage{xcolor}

\newtheorem{theorem}{Theorem}
\newtheorem{proposition}[theorem]{Proposition}
\newtheorem{corollary}[theorem]{Corollary}
\newtheorem{lemma}[theorem]{Lemma}
\newtheorem{definition}[theorem]{Definition}
\theoremstyle{remark}
\newtheorem{remark}[theorem]{Remark}

\title{Note~3: Group Conditional Coverage\\
  in the Non-Gaussian Experiment\\[4pt]
  {\large CKME vs DCP-DR vs hetGP across x-bins}}
\author{Working Note --- CKME-CP Project}
\date{April 2026}

\begin{document}
\maketitle

\begin{abstract}
We analyze \emph{group conditional coverage} --- coverage computed
within equal-width bins of the input space $\mathcal{X} = [0, 2\pi]$ ---
for CKME-CP, DCP-DR, and hetGP across six non-Gaussian simulators.
Using 11 macroreplications with $K=10$ x-bins, we find:
(1)~CKME achieves the most uniform bin-coverage (max bin deviation
$\le 0.10$ across all simulators);
(2)~DCP-DR exhibits severe spatial non-uniformity, with max bin
deviations up to 0.33;
(3)~hetGP over-covers globally on the mixture simulator C1-S
(97.5\% vs 90\% target) but is spatially uniform \emph{conditional on
its elevated level}.
These results support the claim that CKME's CDF-first estimation
produces more spatially homogeneous conformal scores than
quantile-regression-based methods.
\end{abstract}

\tableofcontents
\newpage


%% ====================================================================
\section{Setup}
\label{sec:setup}
%% ====================================================================

\subsection{Simulators}

All six simulators share the same mean function
$f(x) = e^{x/10}\sin(x)$ and target noise scale
$\sigma_{\mathrm{tar}}(x) = 0.1 + 0.1(x-\pi)^2$, but differ in
noise shape.

\medskip
\begin{center}
\begin{tabular}{llll}
\toprule
Name & Noise family & Small (S) parameter & Large (L) parameter \\
\midrule
A1 & Student-$t$ & $\nu=10$ & $\nu=3$ \\
B2 & Centered Gamma & $k=9$ & $k=2$ \\
C1 & Gaussian mixture & $\pi=0.02$ & $\pi=0.10$ \\
\bottomrule
\end{tabular}
\end{center}

\medskip
\noindent
The ``Small'' variants represent mild departures from Gaussianity;
the ``Large'' variants represent strong departures (heavy tails,
strong skewness, or prominent bimodality).

\subsection{Methods}

Three prediction interval methods, all targeting 90\% coverage
($\alpha=0.1$):

\begin{itemize}
  \item \textbf{CKME-CP}: CDF-first conformal prediction with
    abs-median score $s(x,y) = |\hat F(y|x) - 0.5|$.
    Coverage column: \texttt{covered\_score}.
  \item \textbf{DCP-DR}: Distributional conformal prediction with
    doubly-robust quantile regression (R implementation).
    Coverage column: \texttt{covered\_score\_dr}.
  \item \textbf{hetGP}: Heteroscedastic Gaussian process with
    interval-based coverage (R implementation).
    Coverage column: \texttt{covered\_interval\_hetgp}.
\end{itemize}

\subsection{Binning Procedure}

\begin{enumerate}
  \item Partition $[0, 2\pi]$ into $K=10$ equal-width bins:
    $B_k = [\tfrac{2\pi(k-1)}{K},\, \tfrac{2\pi k}{K})$,
    $k=1,\ldots,K$.
  \item For each macrorep $m$ and method, compute bin-coverage:
    \[
      \widehat{\mathrm{Cov}}_k^{(m)}
      = \frac{1}{|I_k^{(m)}|}
        \sum_{i \in I_k^{(m)}} \mathbf{1}\{Y_i \in C(X_i)\},
    \]
    where $I_k^{(m)} = \{i : X_i \in B_k\}$.
  \item Aggregate across $M=11$ macroreps:
    mean $\pm$ std of $\{\widehat{\mathrm{Cov}}_k^{(m)}\}_{m=1}^{M}$.
\end{enumerate}

Each macrorep has $n_{\mathrm{test}} = 1000$ test points with LHS
design on $[0, 2\pi]$, giving approximately 100 points per bin per
macrorep.

\subsection{Metrics}

For each method and simulator, we report:
\begin{itemize}
  \item \textbf{Overall coverage}: mean of bin-coverage across all
    bins (should be near $1-\alpha = 0.90$).
  \item \textbf{Max bin deviation}: $\max_k |\widehat{\mathrm{Cov}}_k
    - (1-\alpha)|$, measuring the worst-case departure from target
    across all bins.
\end{itemize}

A method with good \emph{group conditional coverage} should have both
overall coverage near 90\% and small max bin deviation.


%% ====================================================================
\section{Results}
\label{sec:results}
%% ====================================================================

\subsection{Figure: Bin-Coverage Profiles}

\begin{figure}[htbp]
\centering
\includegraphics[width=\textwidth]{../exp_nongauss/output/group_coverage.png}
\caption{Group conditional coverage by x-bin ($K=10$, 11 macroreps).
  Solid/dashed/dot-dashed lines: mean bin-coverage across macroreps.
  Shaded bands: $\pm 1$ std.
  Black dotted line: 90\% target.
  Top row: Small departure; bottom row: Large departure.}
\label{fig:group-coverage}
\end{figure}

\subsection{Summary Table}

\begin{table}[htbp]
\centering
\caption{Overall coverage and max bin deviation from 90\% target.}
\label{tab:summary}
\begin{tabular}{ll cc}
\toprule
Simulator & Method & Overall Cov.\ (\%) & Max Bin Dev. \\
\midrule
\multicolumn{4}{l}{\emph{Small departure}} \\
A1-S (Student-$t$, $\nu$=10)
  & CKME   & 90.7 & 0.095 \\
  & DCP-DR & 90.9 & 0.212 \\
  & hetGP  & 91.0 & 0.032 \\
\addlinespace
B2-S (Gamma, $k$=9)
  & CKME   & 89.9 & 0.097 \\
  & DCP-DR & 90.1 & 0.245 \\
  & hetGP  & 91.3 & 0.040 \\
\addlinespace
C1-S (Mixture, $\pi$=0.02)
  & CKME   & 89.7 & 0.092 \\
  & DCP-DR & 90.3 & 0.263 \\
  & hetGP  & \textbf{97.5} & 0.080 \\
\midrule
\multicolumn{4}{l}{\emph{Large departure}} \\
A1-L (Student-$t$, $\nu$=3)
  & CKME   & 90.4 & 0.088 \\
  & DCP-DR & 90.4 & 0.189 \\
  & hetGP  & 93.6 & 0.059 \\
\addlinespace
B2-L (Gamma, $k$=2)
  & CKME   & 90.2 & 0.088 \\
  & DCP-DR & 89.9 & 0.297 \\
  & hetGP  & 93.1 & 0.043 \\
\addlinespace
C1-L (Mixture, $\pi$=0.10)
  & CKME   & 89.8 & 0.066 \\
  & DCP-DR & 90.2 & 0.334 \\
  & hetGP  & 90.9 & 0.022 \\
\bottomrule
\end{tabular}
\end{table}


%% ====================================================================
\section{Observations}
\label{sec:observations}
%% ====================================================================

\subsection{CKME: Moderate but Consistent Spatial Uniformity}

Across all six simulators, CKME maintains overall coverage between
89.7\% and 90.7\%, close to the 90\% target. The max bin deviation
ranges from 0.066 (C1-L) to 0.097 (B2-S). The coverage profile in
Figure~\ref{fig:group-coverage} shows a mild ``U-shape'' (slightly
lower coverage near $x = \pi$ where $\sigma(x)$ is smallest), but
the deviations are moderate.

This U-shape is consistent with the fixed-bandwidth effect identified
in Note~2 (Score Homogeneity): at low-$\sigma(x)$ locations, the
effective bandwidth $h_{\mathrm{eff}}(x) = h/\sigma(x)$ is large,
producing a more biased CDF estimate. The abs-median score then
under-covers at these locations.

\subsection{DCP-DR: Severe Spatial Non-Uniformity}

DCP-DR achieves near-target overall coverage (89.9\%--90.9\%) but
exhibits severe spatial variation. Max bin deviations range from
0.189 (A1-L) to \textbf{0.334} (C1-L). The coverage profile shows
a pronounced inverted-U shape: over-coverage near $x=\pi$ and
under-coverage near the domain boundaries.

This pattern is strikingly different from CKME and suggests that the
DCP-DR quantile regression model has spatially varying estimation
accuracy. Near the center of the domain where data density is highest,
the quantile estimates are more accurate and conformal calibration
produces wider-than-necessary intervals (over-coverage). At the
boundaries, quantile extrapolation degrades and coverage drops.

\subsection{hetGP: Spatially Flat but Globally Miscalibrated}

hetGP shows the smallest max bin deviations in most cases (0.022--0.080),
meaning its coverage is spatially \emph{flat}. However, this flatness
comes at the cost of global miscalibration:

\begin{itemize}
  \item On \textbf{C1-S} (mild Gaussian mixture), hetGP dramatically
    over-covers at \textbf{97.5\%} vs the 90\% target. This is because
    hetGP models the noise as Gaussian, and a Gaussian fit to a
    bimodal mixture inflates the estimated variance.
  \item On \textbf{A1-L} and \textbf{B2-L} (heavy-tailed/skewed),
    hetGP over-covers at 93--94\%, again due to Gaussian variance
    inflation when the true distribution has heavier tails.
\end{itemize}

The spatial flatness of hetGP is expected: under a Gaussian model,
the interval width is $\propto \hat\sigma(x)$, which tracks the
noise scale uniformly. The problem is not spatial heterogeneity but
global bias from distributional misspecification.

\subsection{Small vs Large Departure}

Increasing the departure from Gaussianity:
\begin{itemize}
  \item CKME remains stable (max deviation changes by $< 0.03$).
  \item DCP-DR's spatial non-uniformity worsens
    (A1: $0.21 \to 0.19$; B2: $0.25 \to 0.30$; C1: $0.26 \to 0.33$).
  \item hetGP's overall over-coverage increases on A1 and B2
    (91\% $\to$ 93--94\%) but interestingly improves on C1
    (97.5\% $\to$ 90.9\%), likely because the larger mixture
    proportion creates a distribution that is \emph{more} Gaussian-like
    in the aggregate.
\end{itemize}


%% ====================================================================
\section{Group-Level Theory}
\label{sec:group-theory}
%% ====================================================================

Note~2 proves a \emph{pointwise} two-sided conditional coverage bound
(Thm.~2 there) under the adaptive bandwidth $h(x) = c\,\sigma(x)$.
We now extend that result to the group level and give finite-sample
counterparts that match the empirical metric
$\max_k |\widehat{\mathrm{Cov}}_k - (1-\alpha)|$ used in
Section~\ref{sec:results}. All results below assume the location-scale
model of Note~2, Section~2, and write
$\delta(c,n) = \tfrac{\pi^2 c^2}{6}\|g''\|_1 + \|\hat F - F\|_\infty$
for the score-homogeneity gap from~Note~2, Eq.~(5).

% ------------------------------------------------------------------
\subsection{T-G1: Partition-Free Group Coverage}
\label{sec:tg1}
% ------------------------------------------------------------------

\begin{corollary}[Group coverage for any partition]
\label{cor:tg1}
Let $\{B_k\}_{k=1}^{K}$ be any measurable partition of $\mathcal X$
with $\mathbb P(X \in B_k) > 0$. Under the assumptions of Note~2,
Thm.~2,
\begin{equation}
\label{eq:tg1}
  \bigl|\mathbb P(Y \in C(X) \mid X \in B_k)-(1-\alpha)\bigr|
  \;\le\; \delta(c,n) + p_{n+1},
  \qquad k = 1,\ldots,K.
\end{equation}
\end{corollary}

\begin{proof}[Proof]
By Note~2, Thm.~2, $\bigl|\mathbb P(Y\in C(X)\mid X=x) - (1-\alpha)\bigr|
\le \delta(c,n) + p_{n+1}$ for $\mathbb P_X$-almost every $x$. Taking
conditional expectation given $X\in B_k$ and applying the triangle
inequality preserves the same bound.
\end{proof}

\begin{remark}
The bound~\eqref{eq:tg1} holds for \emph{any} partition --- equal-width
bins, quantile bins, or a data-driven slicing --- without retuning.
This is stronger than partition-specific group-CP variants, which must
be re-calibrated per partition.
\end{remark}

% ------------------------------------------------------------------
\subsection{T-G2: Finite-Sample Worst-Bin Bound}
\label{sec:tg2}
% ------------------------------------------------------------------

Corollary~\ref{cor:tg1} controls the \emph{population} bin coverage,
but Table~\ref{tab:summary} reports \emph{empirical} bin coverage
$\widehat{\mathrm{Cov}}_k$ computed from finite test data. The sampling
noise adds a second term.

\begin{theorem}[Worst-bin deviation, finite sample]
\label{thm:tg2}
Let $n_k = |\{i \le n_{\mathrm{test}} : X_i \in B_k\}|$ and
$n_{\min} = \min_k n_k$. Under the assumptions of
Corollary~\ref{cor:tg1}, for any $\eta \in (0,1)$, with probability at
least $1-\eta$,
\begin{equation}
\label{eq:tg2}
  \max_{k \le K}\bigl|\widehat{\mathrm{Cov}}_k - (1-\alpha)\bigr|
  \;\le\;
  \underbrace{\delta(c,n) + p_{n+1}}_{\text{structural (T-G1)}}
  \;+\;
  \underbrace{\sqrt{\tfrac{\log(2K/\eta)}{2\,n_{\min}}}}_{\text{binomial fluctuation}}.
\end{equation}
\end{theorem}

\begin{proof}[Proof sketch]
For each bin $k$, $\widehat{\mathrm{Cov}}_k$ is an empirical mean of
bounded indicators, so Hoeffding's inequality gives
$\mathbb P(|\widehat{\mathrm{Cov}}_k - \mathrm{Cov}_k| \ge t) \le
2\exp(-2 n_k t^2)$. A union bound over $k = 1,\ldots,K$ and
$t = \sqrt{\log(2K/\eta)/(2 n_{\min})}$ gives uniform control.
Combining with Corollary~\ref{cor:tg1} via the triangle inequality
yields~\eqref{eq:tg2}.
\end{proof}

\begin{remark}[Theory--experiment account matching]
\label{rem:account-match}
Theorem~\ref{thm:tg2} decomposes the measured quantity
$\max_k |\widehat{\mathrm{Cov}}_k - (1-\alpha)|$ into a
\emph{structural} term $\delta(c,n) + p_{n+1}$ that the method
controls, and a \emph{sampling} term
$\sqrt{\log(2K/\eta) / (2 n_{\min})}$ that is universal. For the
experimental configuration of Section~\ref{sec:setup}
(\,$K = 10$, $n_k \approx 100$, $\eta = 0.05$), the sampling term
evaluates to
\[
  \sqrt{\tfrac{\log(400)}{200}} \;\approx\; 0.17,
\]
and the empirical CKME max deviations $\le 0.10$ reported in
Table~\ref{tab:summary} are \emph{below} this binomial floor --- i.e.,
structurally indistinguishable from sampling noise. In contrast, the
DCP-DR deviations up to $0.33$ \emph{exceed} the binomial floor by a
factor $\sim 2$, which must be attributed to a non-vanishing
structural term.
\end{remark}

% ------------------------------------------------------------------
\subsection{T-G3: Lower Bound for Quantile-First Scores}
\label{sec:tg3}
% ------------------------------------------------------------------

Remark~\ref{rem:account-match} raises the question of \emph{why}
certain conformal methods exhibit a non-vanishing structural term.
We first analyze the CQR (Conformal Quantized Regression) score
$s_{\mathrm{CQR}}$, which operates directly in Y-space and is the
prototypical quantile-first score. We then discuss how DCP-DR,
despite using a CDF-space score, can suffer from a different source of
non-uniformity.

\begin{proposition}[Y-space scores have a non-vanishing homogeneity gap]
\label{prop:tg3}
Consider the Y-space conformal score
$s(x,y) = \max\bigl(q_{\mathrm{lo}}(x) - y,\;
y - q_{\mathrm{hi}}(x)\bigr)$,
where $q_{\mathrm{lo/hi}}$ are the oracle $\alpha/2$ and
$1-\alpha/2$ quantiles of $Y\mid X$. Under the location-scale model
$Y = f(x) + \sigma(x)\varepsilon$ with $\varepsilon \sim G$,
\begin{equation}
\label{eq:tg3}
  d_{\mathrm{TV}}\bigl(F_{s\mid x},\, F_{s\mid x'}\bigr)
  \;=\; d_{\mathrm{TV}}\bigl(S,\,\rho S\bigr),
  \quad\text{where } \rho = \sigma(x')/\sigma(x),
\end{equation}
and $S := \max\bigl(G^{-1}(\alpha/2) - \varepsilon,\;
\varepsilon - G^{-1}(1-\alpha/2)\bigr)$. In particular
$d_{\mathrm{TV}}(F_{s\mid x}, F_{s\mid x'}) > 0$ whenever
$\sigma(x) \ne \sigma(x')$, and the bound is \emph{independent of $n$}.
\end{proposition}

\begin{proof}[Proof sketch]
Substituting the oracle quantiles and factoring $\sigma(x)$:
$s(x, Y) = \sigma(x)\,\max\bigl(G^{-1}(\alpha/2) - \varepsilon,\;
\varepsilon - G^{-1}(1-\alpha/2)\bigr) = \sigma(x)\,S$. Hence
$F_{s\mid x}$ is the distribution of $\sigma(x)\,S$, and
$d_{\mathrm{TV}}(\sigma S, \sigma' S) = d_{\mathrm{TV}}(S, \rho S)$.
The TV distance between $S$ and a scaled version is strictly positive
unless $\rho = 1$. Since the derivation uses only the oracle
quantiles, no amount of data can shrink this gap.
\end{proof}

\begin{remark}[Numerical illustration]
\label{rem:tg3-numerical}
In our simulators,
$\sigma_{\mathrm{tar}}(x) = 0.1 + 0.1(x-\pi)^2$ varies between $0.1$
(at $x=\pi$) and $\approx 1.09$ (at $x=0$ or $2\pi$), giving
$\rho_{\max} \approx 10.9$. A direct calculation for Gaussian
$\varepsilon$ and $\alpha = 0.1$ gives
$d_{\mathrm{TV}}(S, 10.9\,S) \approx 0.45$. Thus any CQR-style
conformal method would face a worst-case group coverage deviation of
at least $0.45$ on these simulators, regardless of sample size.
\end{remark}

\begin{remark}[CDF-space scores and model misspecification]
\label{rem:tg3-dcpdr}
Proposition~\ref{prop:tg3} applies to Y-space scores such as CQR.
DCP-DR uses a \emph{different} score: $s(x,y) = |\hat F(y|x) - 0.5|$,
which is a CDF-space score identical in form to CKME's abs-median
score. Under the oracle CDF, the PIT guarantees
$F(Y|X=x) \sim \mathrm{Uniform}(0,1)$, so this score would be
perfectly homogeneous.

However, DCP-DR estimates $\hat F$ via linear logistic regression:
$\mathrm{logit}(\hat F(y|x)) = \beta_0(y) + \beta_1(y)\,x$.
When the true conditional CDF is nonlinear in $x$, this linear model
is misspecified and the approximation error
$\|\hat F(\cdot|x) - F(\cdot|x)\|_\infty$ varies across $x$. The
resulting score heterogeneity is an \emph{estimation gap} --- in
principle eliminable with a more flexible CDF estimator --- rather
than the \emph{structural gap} of Proposition~\ref{prop:tg3}, which
persists even under oracle quantiles. The empirical DCP-DR max
deviation of $0.33$ in our experiments reflects this model
misspecification effect.
\end{remark}

% ------------------------------------------------------------------
\subsection{T-G4: Connection to Distribution-Free Impossibility}
\label{sec:tg4}
% ------------------------------------------------------------------

\begin{remark}[Structural assumptions as an escape hatch]
\label{rem:tg4}
Barber, Cauchois, Romano, and Cand\`es (\emph{Ann.\ Stat.}, 2021) prove that in
the distribution-free setting, non-trivial \emph{pointwise}
conditional coverage is impossible. Group conditional coverage with a
fixed partition $\{B_k\}$ is a tractable relaxation, typically
achieved by running split conformal separately on each bin (``group
conformal'').

Corollary~\ref{cor:tg1} provides a different escape from the
impossibility: rather than partitioning the calibration data, we add
a structural assumption (location-scale noise plus CDF-first scores
with adaptive bandwidth). Under this assumption a single pooled
calibration quantile simultaneously controls coverage on every
partition, with the same structural bound. The price is the
parametric model; the reward is partition-freedom
(Corollary~\ref{cor:tg1}) and, crucially, no split of the
calibration budget across bins.
\end{remark}


%% ====================================================================
\section{Discussion}
\label{sec:discussion}
%% ====================================================================

\subsection{Connection to Theory}

The theoretical framework of Section~\ref{sec:group-theory} lets us
account for the empirical patterns in Section~\ref{sec:results} term
by term.

\paragraph{CKME is at the binomial floor (T-G2).}
Remark~\ref{rem:account-match} shows that with $K=10$ and
$n_k \approx 100$, the binomial fluctuation term in
Theorem~\ref{thm:tg2} alone accounts for $\approx 0.17$ of worst-bin
deviation. CKME's empirical $\max_k |\widehat{\mathrm{Cov}}_k -
(1-\alpha)| \le 0.10$ lies entirely below this floor, meaning the
\emph{structural} term $\delta(c,n) + p_{n+1}$ is already close to
zero on all six simulators.

\paragraph{DCP-DR carries an estimation gap (T-G3, Remark~\ref{rem:tg3-dcpdr}).}
DCP-DR uses the same CDF-space score $|\hat F - 0.5|$ as CKME, so the
\emph{structural} impossibility of Proposition~\ref{prop:tg3} (which
targets Y-space CQR scores) does not directly apply. Instead,
Remark~\ref{rem:tg3-dcpdr} identifies DCP-DR's gap as an
\emph{estimation} artifact: its linear logistic CDF model is
misspecified for the nonlinear $f(x)$, producing spatially varying
$\|\hat F - F\|_\infty$. The empirical max deviation of $0.33$ reflects
this model misspecification, not an intrinsic score-construction
limitation.

\paragraph{hetGP is globally biased but spatially flat.}
hetGP does not fit either framework: its score is model-based
(parametric Gaussian likelihood) rather than conformal, so
Corollary~\ref{cor:tg1} and Proposition~\ref{prop:tg3} do not apply.
The small max bin deviation reflects the fact that under a Gaussian
model the interval width is proportional to $\hat\sigma(x)$
\emph{uniformly}, while the global over-coverage reflects
distributional misspecification.

\paragraph{Caveat on adaptive bandwidth.}
Note that Corollary~\ref{cor:tg1} requires $h(x) = c\,\sigma(x)$,
while the main experiment of Section~\ref{sec:results} uses a
\emph{fixed} CV-tuned bandwidth. That CKME nevertheless achieves
structural term $\approx 0$ suggests either (i) the fixed $h$ lies
close enough to $c\,\sigma(x)$ at the relevant scale, or (ii) bin
averaging in~\eqref{eq:tg1} is slack enough to absorb a moderate
pointwise $d_{\mathrm{TV}}$. A direct adaptive-$h$ ablation on B2L
and C1L is deferred to Section~\ref{sec:discussion}~Limitations.

\subsection{Method Comparison Summary}

\begin{center}
\begin{tabular}{lcc}
\toprule
Method & Overall calibration & Spatial uniformity \\
\midrule
CKME-CP & \checkmark (89.7--90.7\%) & Moderate (max dev $\le$ 0.10) \\
DCP-DR  & \checkmark (89.9--90.9\%) & Poor (max dev up to 0.33) \\
hetGP   & $\times$ (up to 97.5\%) & Good (max dev $\le$ 0.08) \\
\bottomrule
\end{tabular}
\end{center}

\medskip
\noindent
CKME is the only method that achieves \emph{both} good overall
calibration and reasonable spatial uniformity across all noise types.

\subsection{Limitations and Next Steps}

\begin{itemize}
  \item Only 11 macroreps --- larger runs (50+) would tighten the
    std bands and confirm the patterns.
  \item All simulators are 1D ($x \in [0, 2\pi]$). The 2D
    \texttt{exp\_branin} experiment would test whether spatial
    uniformity holds in higher dimensions.
  \item The analysis uses equal-width bins. For non-uniform designs,
    equal-count (quantile) bins may be more appropriate.
  \item Comparison against dedicated group-coverage methods (RLCP,
    conditional-conformal) is planned as future work.
\end{itemize}


\end{document}
